\chapter{From Memory to Pages}
\label{ch:pagestore}

OntoDB now has two complete recovery designs. The memory database of
\cref{ch:durability} is redo-only: no steal, a logical log, and snapshots. The
page store of \cref{ch:aries} steals pages, logs page images, and recovers
with ARIES, but only tests open it. This chapter is about the step that joins
the second design to the shell, so that \code{ontodb <dir>} keeps its data on
pages. Most of the work is not new algorithms; it is deciding what a directory
\emph{is}, what opening one does, and which operations get which kind of log
record. Those decisions are where real systems differ.

\section{What is in a directory}

A database directory must say which design wrote it, because the two have
incompatible files. The memory database keeps \code{snapshot} and \code{wal};
the page store keeps \code{db}, the database file of \cref{ch:pages}, and its
own \code{wal}. A one-line \code{format} file names the store and its version,
\code{pages 1}, much as PostgreSQL keeps a \code{PG\_VERSION} file in every
data directory.

\begin{center}
\begin{tabular}{@{}ll@{}}
\toprule
Directory holds & Opens as \\
\midrule
\code{format} = \code{pages 1} & a page store \\
\code{snapshot}, no \code{format} & a memory database, as today \\
\code{format} naming another store or version & refused, saying which \\
nothing & a new page store (or memory, with \code{--memory}) \\
\bottomrule
\end{tabular}
\end{center}

Refusing an unknown format is the important row. A newer OntoDB that changes
the page layout must not be opened by an older one that would misread it, and
the cheapest protection is a version number checked before any page is read.

\begin{decision}{A format file names the store (planned)}
\dpart{Choice} New directories get \code{format} with the store and version;
a directory without it but with a \code{snapshot} is a memory database.
\dpart{Alternatives} Detect the store from which files exist; or put the
version in the database file's header page.
\dpart{Why} Detection by file names breaks the day both designs share a name,
and a header page cannot describe a store that has no pages. A separate file
is read before anything else and is trivial to inspect.
\dpart{Cost} One more file to create atomically with the others.
\end{decision}

\section{Opening is recovery}

On the memory path, opening means loading a snapshot and replaying the log.
On the page path, opening a directory runs all of \cref{ch:aries}, in an order
that matters:

\begin{enumerate}
\item \textbf{Cut a torn tail.} Scan the log and truncate it at the first record
that is incomplete or fails its checksum. The recovery modes of
\cref{ch:durability} apply: point in time truncates and reports; \code{--strict}
refuses.
\item \textbf{Analysis and redo.} Rebuild the transaction table and the dirty
page table, then repeat history from the smallest recLSN, so every page is as
it was at the crash, uncommitted changes included.
\item \textbf{Open the catalog, the dictionary, and the indexes.} Only now are
the pages that describe them trustworthy, since redo has brought them up to
date.
\item \textbf{Undo.} Roll back the loser transactions, writing compensation
records, exactly as an abort does.
\end{enumerate}

Step 3 sits between redo and undo for a reason. Redo is physical: it needs
page ids and bytes, not trees. Undo of an index change is logical, ``remove
this key from that tree'', because a split may have moved the key to another
page since it was inserted; logical undo needs the tree, and the tree needs
its root page, which the catalog names.

\begin{keyidea}
Redo works on pages and needs nothing above the buffer pool. Undo works on
structures and needs the access methods. Recovery opens the engine layer by
layer, in that order.
\end{keyidea}

\section{Statements log on pages}

On a page directory the transaction manager is given the log manager, and from
then on the machinery of \cref{ch:wal,ch:aries} runs for every statement:
each page change logs a record, commit forces the log, and abort undoes by the
log. The dictionary's pages are pages like any other, so a new term is covered
by the page records of the term heap and the dictionary tree; there is no
separate term record as on the memory path. A term added by a transaction that
then aborts stays in the dictionary, for the same reason as there: ids are
handed out densely, and an unused term is invisible.

The test for this step is blunt and sufficient: commit a triple, kill the
process before any page is written, reopen, and find the triple. Only the log
can have put it there.

\section{Loading: one record per triple, or one per load}

\code{\textbackslash load} of a large file is one transaction, and it is where
the two obvious logging designs diverge most:

\begin{description}[leftmargin=1.2em]
\item[Log every change.] Each inserted triple dirties leaf pages of three
trees, and each change is logged. Simple and uniform: a load is just a big
transaction. The log grows with the data, roughly three page records per
triple until pages fill, and the load runs at the speed of logging.
\item[Minimal logging.] Build the trees bottom-up from sorted input
(\cref{ch:btree}'s bulk load), write every page to the file, force the file,
and log one record that names the load. Undo of the load is ``drop the new
trees''. SQL Server calls this minimal logging, and PostgreSQL, with
\code{wal\_level} set to \code{minimal}, skips the log for a table created in
the same transaction, for the same reason.
\end{description}

Minimal logging is only correct when the trees it builds start empty (or are
new trees swapped in at commit): there is nothing older on those pages for a
partial load to have damaged, so undo can discard them whole.

\begin{decision}{How a load is logged (planned, to be measured)}
\dpart{Choice} Decide between per-triple logging and minimal logging by
measuring both on LUBM(1): log bytes and load time. A crash during a load must
leave the database as it was before it, either way.
\dpart{Alternatives} The two above; or per-page images of the newly built
trees.
\dpart{Why} The answer depends on numbers this engine has not produced yet:
how much faster a bulk build is than three tree inserts per triple, and how
much log a load writes.
\dpart{Cost} Minimal logging needs empty target trees and forces every page
before commit; per-triple logging makes the log as large as the data.
\end{decision}

\section{Closing cleanly, and checkpoints}

A clean exit (\code{\textbackslash quit}, or end of input) writes every dirty
page, forces the log, and closes the file. The next open then finds nothing to
redo, and its summary says so; that is the cheapest test that the close was
clean. \code{\textbackslash checkpoint} runs the sharp checkpoint of
\cref{ch:aries} with every statement held out, because a sharp checkpoint must
not see a page change in the middle of being made. \Cref{ch:logkeep} replaces
it with a fuzzy one.

\section{Checking a page database}

\code{\textbackslash verify} keeps its promise of changing nothing and grows
three checks: every tree satisfies its invariants (\cref{ch:btree}); SPO, POS,
and OSP hold the same set of triples; and every id in them is in the
dictionary. The second check is the one specific to RDF: three indexes that
disagree return different answers depending on which one the optimizer
chooses, the worst kind of bug to find from the query side.

\section{Testing through the shell}

The crash harness of \cref{ch:aries} kills a store driven by test code. The
last step drives the real \code{Database} with SPARQL text on a page directory
and kills it between statements and at named crash points. Test-only options
make tree nodes tiny and the buffer pool small, so a short script splits,
merges, and steals pages constantly; otherwise a test database never leaves its
first page, and the hard paths never run.

\begin{status}
Planned, and next (plan 10.1--10.7). Today \code{ontodb <dir>} opens the memory
database of \cref{ch:durability}, and the page store's recovery runs only in
tests and in the crash harness.
\end{status}

\begin{checkyourself}
\item Why must an older OntoDB refuse to open a directory whose \code{format} names a newer version, rather than try?
\item Why are the catalog and the indexes opened after redo but before undo?
\item A transaction commits one \code{INSERT DATA}, and the process is killed before any page is written. Where does the reopened database find the triple?
\item Under what condition is minimal logging of a load correct, and what does undo of such a load do?
\item What does a clean close guarantee about the next open, and how can a test observe it?
\item Give a query whose answer depends on which of SPO, POS, or OSP the optimizer scans if the three disagree.
\end{checkyourself}

\section*{Further reading}
Mohan et al.\ \cite{aries} on logical undo of index changes; Gray and Reuter
\cite{grayreuter} on recovery-time ordering; the SQL Server documentation on
minimally logged operations \cite{sqlminimal}.
